% The next line makes it so it compiles the main file if I hit the compile button
% when editing this file in TeXworks.
% !TEX root = ../ActiveIntroToDiscreteMathAndAlgorithms.tex
%%%%%%%%%%%%%%%%%%%%%%%%%%%%%%%%%%%%
\section{Problems}

\begin{pro}
Use induction to prove that  $\dsum _{k=1}^n k^3 = \dfrac{n^2(n+1)^2}{4}$ 
for all $n\geq 1$.
\end{pro}

\begin{pro}
Use induction to prove that for all $n\geq 2$,
$$\dsum _{k=2}^n \dfrac{1}{(k-1)k}=\dfrac{1}{1\cdot 2} +
\dfrac{1}{2\cdot 3} + \dfrac{1}{3\cdot 4} + \cdots +
\dfrac{1}{(n-1)\cdot n} = \dfrac{n-1}{n}.$$
\end{pro}

%It's an example, so it doesn't make a very good problem.
%\begin{pro}[Towers of Hanoi]
%The following legend is attributed to
% French mathematician Edouard Lucas in 1883. 
%%the toy seen in figure \ref{fig:towers-of-hanoi}  (with eight disks), along with the
%%following  legend. 
%The tower of Brahma  had 64 disks of gold resting on three diamond needles. 
%At the beginning of time, God placed these disks on the first needle
%and ordained that a group of priests should transfer them to the third needle according to the following rules:
%\begin{enumerate}
%\item The disks are initially stacked on peg A, in decreasing order (from bottom to top).
%\item The disks must be moved to another peg in such a way that only one disk is moved at a time and without stacking a larger disk onto a smaller disk.
%\end{enumerate}
%When they finish, the Tower will crumble and the world will end. Prove that if there are $n$ disks, then $2^n-1$ moves 
%are necessary and sufficient to perform the task according to the rules.
%(Hint: First develop a recursive algorithm for the problem.  Then find  a recurrence relation for the number of times a disk is moved
%in your recursive algorithm.  Finally, use induction to prove that the solution to the recurrence relation is $2^n-1$.
%\end{pro}
%\vspace{1cm}
% This is being a pain in the butt, so I removed it.
%\begin{figurehere}
%\centering
%\includegraphics[scale=.2]{tower-of-hanoi2.eps}
%\vspace{1cm} 
%%hangcaption{Towers of Hanoi.} 
%\label{fig:towers-of-hanoi}
%\end{figurehere}
%\vspace{1cm}
%
%Probably too subtle of a problem.
%\begin{pro}[Josephus' Problem] In \cite{HeKa} we find the following legend about the famous first-century Jewish historian
%Flavius Josephus:
%\begin{quote}
%In the Jewish revolt against Rome, Josephus and 39 of his comrades were holding out against the Romans in a cave. With defeat imminent,
%they resolved that, like the rebels at Masada, they would rather die than be slaves to the Romans.
%They decided to arrange themselves in a circle. One man was designated as number one, and they proceeded clockwise killing every
%seventh man\ldots . Josephus (according to the story) was among other things an accomplished mathematician; so he instantly
%figured out where he ought to sit in order to be the last to go. But when the time came, instead of killing himself he joined the Roman side.
%\end{quote}
%In general, given a group of $n$ men arranged in a circle under the edict that every $m$th
%man will be executed going around the circle until only one remains,
%the object is to find the position $L(n,m)$ in which you should stand in order to be the last survivor.
%The particular situation of Flavius Josephus is  asking for $L(40,7)$.
%The general Josephus' Problem is very difficult.
%Prove, however, that $L(n,2)=1+2n-2^{1+\floor{\log_2n}}$.
%\end{pro}

%My problems.


\begin{pro}
Prove that for all positive integers $n$, $f_1^2 + f_2^2 + \cdots +f_n^2 = f_nf_{n+1}$,  
where $f_n$ is the $n$th Fibonacci number.
\end{pro}



\begin{pro}
Prove the following generalized De Morgan's Law for sets (where $n\geq 2$):
$$\overline{(A_1\cup A_2\cup\cdots \cup A_n)} = \overline{A_1}\cap \overline{A_2}\cap\cdots\cap \overline{A_n}.$$
(Note: There is a second law just like it that swaps the $\cap$s and $\cup$s.)
\end{pro}

\begin{pro}
Prove that if $n\geq 4$, $n!> 2^n$.
\end{pro}

\begin{pro}
Prove that the number of binary palindromes of length $2k+1$ (odd length) is 
$2^{k+1}$ for all $k\geq 0$.
\end{pro}

\begin{pro}
Prove that a set of size $n\geq 1$ has $2^n$ subsets.
\end{pro}

%now an example.
%\begin{pro}
%Use induction (on the size of the array) to prove that 
%the {\em recursive} \verb+binarySearch+ algorithm from
%Example~\ref{exa:binSearchComp} is correct.
%(Hint:  If you do not use the fact that the array is sorted, your proof cannot
%possibly be correct.)
%\end{pro}

\begin{pro}
Prove that the \verb+FibR(n)+ algorithm from Example~\ref{exa:fibAlg} correctly computes $f_n$.
(Hint:  Use induction.  How many base cases do you need?  Do you need weak or strong
induction?)
\end{pro}


\begin{pro}
In Example~\ref{exa:fibSol} we gave a solution to the recurrence $f_n=f_{n-1}+f_{n-2}$,
$f_0=0$, $f_1=1$.  Use the substitution method to re-prove this.
(Hint: Recall that  the
roots to the polynomial $x^2-x-1=0$ are $\frac{1 \pm \sqrt{5}}{2}$.  
This is equivalent to $x^2=x+1$. 
You will find this helpful in the inductive step of the proof. 
\end{pro}

\begin{pro}
Explain why the following joke never ends: 
{\em Pete and Repete got in a boat.  Pete fell off.  Who's left?}.
\end{pro}

\begin{pro}
Find and prove a solution for each of the following recurrence relations using two different
techniques (this will not only help you verify that your solutions are correct, but
it will also give you more practice using each of the techniques).  
{\em At least one of the techniques must yield an exact formula if possible.}
\begin{lenum}
\item $T(n)=T(n/2)+n^2$, $T(1)=1$. (You may assume $n$ is a power of 2.)
\item $T(n)=T(n/2)+n$, $T(1)=1$. (You may assume $n$ is a power of 2.)
\item $T(n)=2T(n/2)+n^2$, $T(1)=1$. (You may assume $n$ is a power of 2.)
\item $T(n)=T(n-1)\cdot T(n-2)$, $T(0)=1$, $T(1)=2$. %T(n)=2^{f_n}.
\item $T(n)=T(n-1)+n^2$, $T(1)=1$. % Sum of the first $n$ squares.
\item $T(n)=T(n-1)+2n$, $T(1)=2$.
\end{lenum}

\end{pro}


\begin{pro}
Give an exact solution for each of the following recurrence relations.
\begin{lenum}
\item $a_n=3a_{n-1}$, $a_1=5$.
\item $a_n=3a_{n-1}+2n$, $a_1=5$.
\item $a_n=a_{n-1}+2a_{n-2}$, $a_0=2$, $a_1=5$.
\item $a_n=6a_{n-1}+9a_{n-2}$, $a_0=1$, $a_1=2$.
\item $a_n=-a_{n-1}+6a_{n-2}$, $a_0=4$, $a_1=5$.
\end{lenum}

\end{pro}

\begin{pro}
Use the Master Theorem to find a tight bound for each of the following
recurrence relations.
\begin{lenum}
\item $T(n)=8T(n/2)+7n^3+6n^2+5n+4$.
\item $T(n)=3T(n/5)+n^2-4n+23$.
\item $T(n)=3T(n/2)+3$.
\item $T(n)=T(n/3)+n$.
\item $T(n)=2T(2n/5)+n$.
\item $T(n)=5T(2n/5)+n$.
\end{lenum}

\end{pro}

\begin{pro}\label{pro:rpartComp}
Prove that the \scode{Partition} algorithm from Example~\ref{exa:quicksortImp} has complexity $\Theta(n)$.
\end{pro}

\begin{pro}
Consider the classic {\em bubble sort} algorithm (see Example~\ref{exa:bubble}).
\begin{lenum}
\item
Write a recursive version of the bubble sort algorithm.
(Hint:  The algorithm I have in mind should contain one recursive call and one loop.)
\item
Let $B(n)$ be the complexity of your recursive version of bubble sort.
Give a recurrence relation for $B(n)$.
\item
Solve the recurrence relation for $B(n)$ that you developed in part (b).
\item
Is your recursive implementation better, worse, or the same as the iterative one given
in Example~\ref{exa:bubble}?  Justify your answer.
\end{lenum}
\end{pro}

\begin{pro}
Consider the following algorithm (remember that integer division truncates):
\begin{code}
    int halfIt(int n) {
        if(n>0) {
            return 1 + halfIt(n/2);
        } else {
            return 0;
        }
    }
\end{code}
\begin{lenum}
\item
What does \verb+halfIt(n)+ return?  Your answer should be a function of $n$.
\item
Prove that the algorithm is correct.  That is, prove that it returns the
answer you gave in part (a).
\item
What is the complexity of \verb+halfIt(n)+? Give and prove an exact formula.
(Hint: This will probably involve developing and solving a recurrence relation.)
\end{lenum}
\end{pro}

\begin{pro}
This problem involves an algorithm to compute the sum of the first $n$ squares
$\left(\mbox{i.e. }\dsum_{k=1}^n k^2\right)$ using recursion.
\begin{lenum}
\item
Write an algorithm to compute $\dsum_{k=1}^n k^2$ that uses recursion and
only uses the increment/decrement operator for arithmetic (e.g., you cannot use addition or multiplication).
(Hint:  The algorithm I have in mind has one recursive call and one or two loops. You should NOT need to 
pass in multiple variables, and you should NOT need a global variable.)
\item
Let $S(n)$ be the complexity of your algorithm from part (a).
Give a recurrence relation for $S(n)$.
\item
Solve the recurrence relation for $S(n)$ that you developed in part (b) to determine the complexity 
of your algorithm in part (a).
\item
Give a recursive linear-time algorithm to solve this same problem
(with no restrictions on what operations you may use).
Prove that the algorithm is linear.
\item
Give a constant-time algorithm to solve this same problem
(with no restrictions on what you may use).
Prove that the algorithm is constant.
\item
Discuss the relative merits of the three algorithms.  Which algorithm is best?
Worst?  Justify.
\end{lenum}
\end{pro}

\begin{pro}
Assuming the priests can move one disk per second, that they started moving disks
6000 years ago, and that the legend of the Towers of Hanoi is true, 
when will the world end?
\end{pro}

\begin{pro}
Prove that the \verb+stoogeSort+ algorithm given in
Exercise~\ref{exer:stooge} correctly sorts an array of $n$ integers.
\end{pro}



\begin{pro}
Prove, using mathematical induction, that $\begin{bmatrix}1 & 1
\cr 0 & 1 \cr
\end{bmatrix}^n=\begin{bmatrix}1 & n \cr 0 & 1 \cr
\end{bmatrix}$.
\end{pro}

\begin{pro}
Let $A = \begin{bmatrix} 0 & 3 \cr 2 & 0 \cr \end{bmatrix}$. Find,
with proof, $A^{2003}$.
%\begin{answer}
%$A^{2003} = \begin{bmatrix} 0 & 2^{1001}3^{1002} \cr
%2^{1002}3^{1001} & 0 \cr
%\end{bmatrix}$.
%\end{answer}
\end{pro}

\begin{pro}
(Requires some trig identities you may have forgotten!)
 Let $A=
\begin{bmatrix} \cos \alpha & -\sin \alpha \cr \sin \alpha & \cos \alpha\end{bmatrix}.$
Prove that for $n \in \BBN , n \geq 1$,
$A^n = \begin{bmatrix} \cos n\alpha & -\sin n\alpha \cr \sin n\alpha & \cos n\alpha \cr \end{bmatrix}.$
(Do I have to give you a hint to use induction?)
%\begin{answer}The assertion is clearly true for $n = 1.$
%Assume that is it true for $n$, that is, assume
%$$A^{n} = \begin{bmatrix} \cos (n)\alpha & -\sin (n)\alpha \cr \sin (n)\alpha & \cos (n)\alpha \cr \end{bmatrix}.$$
%Then $$\begin{array}{lll} A^{n+1} & = & AA^{n} \\
%& = & \begin{bmatrix} \cos \alpha & -\sin \alpha \cr \sin \alpha &
%\cos \alpha \cr \end{bmatrix}\begin{bmatrix} \cos (n)\alpha &
%-\sin (n)\alpha \cr \sin (n)\alpha & \cos (n)\alpha \cr
%\end{bmatrix}\vspace{2mm} \\
%& = & \begin{bmatrix} \cos\alpha\cos (n)\alpha  - \sin\alpha\sin
%(n)\alpha& -\cos\alpha\sin (n)\alpha  - \sin\alpha\cos
%(n)\alpha\cr \sin\alpha\cos (n)\alpha  + \cos\alpha\sin (n)\alpha&
%-\sin\alpha\sin (n)\alpha + \cos\alpha\cos (n)\alpha \cr
%\end{bmatrix} \vspace{2mm} \\
%& = & \begin{bmatrix} \cos (n+1)\alpha & -\sin (n+1)\alpha \cr
%\sin (n+1)\alpha & \cos (n+1)\alpha \cr \end{bmatrix},
%\end{array}$$and the result follows by induction.
%\end{answer}
\end{pro}

\begin{pro}
Let $A,B\in\mats{n\times n}$ and $k$ be a positive
integer such that $A^k = {\bf 0}_n$. 
Prove that if $AB = B$ then $B ={\bf 0}_n$.
(Hint: Try induction.)
%\begin{answer}
%Argue inductively, $$\begin{array}{l} A^2B = A(AB) = AB = B \\
%A^3B = A(A^2B) = A(AB) = AB = B \\ \vdots \\ A^mB = AB = B.
%\end{array}$$Hence $B = A^mB = {\bf 0}_nB = {\bf 0}_n$.
%\end{answer}
\end{pro}

\begin{pro}
Prove, by means of induction, that for the following $n\times n$
matrix we have
$$\begin{bmatrix} 1 & 1 & 1 & \cdots & 1 \cr 0 & 1 & 1 & \cdots & 1 \cr
0 & 0 & 1 & \cdots & 1 \cr \vdots & \vdots & \vdots & \cdots &
\vdots \cr 0 & 0 & 0 & \cdots & 1 \cr
\end{bmatrix}^3 = \begin{bmatrix}
1 & 3 & 6 & \cdots & \frac{n(n+1)}{2} \cr 0 & 1 & 3 & \cdots &
\frac{(n-1)n}{2} \cr 0 & 0 & 1 & \cdots & \frac{(n-2)(n-1)}{2} \cr
\vdots & \vdots & \cdots & \cdots & \vdots \cr 0 & 0 & 0 & \cdots &
1 \cr\end{bmatrix}.
$$
\end{pro}


\begin{pro}
Let $$A = \begin{bmatrix} 1 & -1 & -1 \cr -1 & 1 & -1 \cr -1 & -1 &
1 \cr\end{bmatrix}.$$
Conjecture a formula for $A^n$ and prove it
using induction.
%\begin{answer}
%Observe that $A = 2{\bf I}_3 - J$, where ${\bf I}_3$  is the
%$3\times 3$ identity matrix and $$J = \begin{bmatrix} 1 & 1 & 1 \cr
%1 & 1 & 1 \cr 1 & 1 & 1 \cr\end{bmatrix}.$$ Observe that $J^k =
%3^{k-1}J$ for integer $k \geq 1$. Using the binomial theorem we have
%$$\begin{array}{lll}A^n &= & (2{\bf I}_3 - J)^n \\
%&= & \sum _{k=0} ^n \binom{n}{k} (2{\bf I}_3)^{n-k}(-1)^kJ^k \\
%& = & 2^n{\bf I}_3+\dfrac{1}{3}J\sum _{k=1} ^n \binom{n}{k}
%2^{n-k}(-1)^k3^{k} \\
%& = & 2^n{\bf I}_3+ \dfrac{1}{3}J((-1)^n-2^n) \\
% & = & \dfrac{1}{3}\begin{bmatrix} (-1)^n+2^{n+1} & (-1)^n-2^n & (-1)^n-2^n \cr
%(-1)^n-2^n& (-1)^n+2^{n+1} & (-1)^n-2^n \cr (-1)^n-2^n & (-1)^n-2^n
%& (-1)^n+2^{n+1} \cr\end{bmatrix}.\end{array}$$
%\end{answer}
\end{pro}

\begin{pro}
Recall Procedure~\ref{proc:euclid} that gave Euclid's algorithm.
Give a recursive implementation of that algorithm. 
\end{pro}
